\documentclass[12pt,a4paper]{article}
\usepackage[UTF8]{ctex}
\usepackage{amsmath,amssymb,amsthm}
\usepackage{geometry}
\usepackage{bm}
\usepackage{hyperref}

\geometry{margin=2.5cm}

\title{例12.11\quad 装配的稳定性\quad 详解}
\author{现代机器人学：第12章 抓取与操作\\
详细解释}
\date{}

\begin{document}

\maketitle

\section{问题描述}

考虑图12.27中的拱结构。问题是：\textbf{它在重力下稳定吗？}

\subsection{系统组成}

拱结构由多个物体组成：
\begin{itemize}
    \item 物体1：拱的左侧部分
    \item 物体2：拱的右侧部分
    \item 物体3：拱的顶部（可能还有其他物体）
    \item 物体之间存在接触
    \item 每个物体都受到重力作用
\end{itemize}

\section{问题特点}

\subsection{多物体系统}

这是一个\textbf{多物体系统}，与前面的单物体例子不同：
\begin{itemize}
    \item 有多个可移动的物体
    \item 物体之间存在接触约束
    \item 每个物体都有自己的动力学方程
    \item 接触力是成对出现的（作用-反作用）
\end{itemize}

\subsection{图形方法的局限性}

对于这样的问题，图形平面方法很难使用：
\begin{itemize}
    \item 有多个移动物体
    \item 接触关系复杂
    \item 难以在平面上可视化所有约束
\end{itemize}

因此，我们需要使用\textbf{代数方法}。

\section{分析方法}

\subsection{接触模式假设}

我们代数测试所有接触标记为$R$（滚动/固定）的接触模式的一致性。

\begin{itemize}
    \item 假设所有接触都是固定的（$R$模式）
    \item 测试是否存在满足所有约束的一致解
    \item 如果存在，则装配是稳定的
    \item 如果不存在，则装配不稳定
\end{itemize}

\subsection{摩擦锥}

摩擦锥如图12.27所示。每个接触都有摩擦，摩擦锥定义了可以通过该接触传递的力的范围。

\section{数学建模}

\subsection{接触编号}

假设总共有16个接触（或接触wrench）：
\begin{itemize}
    \item 接触1-8：与物体1相关的接触
    \item 接触9-16：与物体2相关的接触
    \item 接触5-12：物体1和物体3之间的接触，以及物体2和物体3之间的接触
\end{itemize}

\subsection{Wrench平衡方程}

对于每个物体，需要满足wrench平衡。假设所有接触都是$R$模式（固定接触），则：

\subsubsection{物体1的平衡}

\begin{equation}
\sum_{i=1}^8 k_i F_i + F_{\text{ext}1} = 0
\label{eq:body1}
\end{equation}

其中：
\begin{itemize}
    \item $k_i \geq 0$，$i = 1, \ldots, 8$是接触系数
    \item $F_i$是接触$i$的wrench（在摩擦锥内）
    \item $F_{\text{ext}1}$是物体1上的重力wrench
\end{itemize}

\subsubsection{物体2的平衡}

\begin{equation}
\sum_{i=9}^{16} k_i F_i + F_{\text{ext}2} = 0
\label{eq:body2}
\end{equation}

其中：
\begin{itemize}
    \item $k_i \geq 0$，$i = 9, \ldots, 16$是接触系数
    \item $F_i$是接触$i$的wrench（在摩擦锥内）
    \item $F_{\text{ext}2}$是物体2上的重力wrench
\end{itemize}

\subsubsection{物体3的平衡}

\begin{equation}
-\sum_{i=5}^{12} k_i F_i + F_{\text{ext}3} = 0
\label{eq:body3}
\end{equation}

其中：
\begin{itemize}
    \item 负号是因为物体3施加到物体1和2的wrench与物体1和2施加到物体3的wrench方向相反
    \item $F_{\text{ext}3}$是物体3上的重力wrench
\end{itemize}

\subsection{作用-反作用原理}

最后一组方程来自\textbf{作用-反作用原理}：

\begin{itemize}
    \item 物体1施加到物体3的wrench等于且相反于物体3施加到物体1的wrench
    \item 物体2施加到物体3的wrench等于且相反于物体3施加到物体2的wrench
    \item 因此，物体3的平衡方程中使用负号
\end{itemize}

\section{完整的问题表述}

\subsection{线性约束满足问题}

我们需要找到$k_i \geq 0$，$i = 1, \ldots, 16$，满足：

\begin{align}
\sum_{i=1}^8 k_i F_i + F_{\text{ext}1} &= 0 \label{eq:system1} \\
\sum_{i=9}^{16} k_i F_i + F_{\text{ext}2} &= 0 \label{eq:system2} \\
-\sum_{i=5}^{12} k_i F_i + F_{\text{ext}3} &= 0 \label{eq:system3}
\end{align}

其中：
\begin{itemize}
    \item 每个方程是3维的（平面问题）或6维的（三维问题）
    \item 总共有9个标量方程（3个物体 $\times$ 3个分量）或18个标量方程（3个物体 $\times$ 6个分量）
    \item 有16个未知数$k_i$
    \item 约束条件：$k_i \geq 0$，$F_i \in \mathcal{WC}_i$（在摩擦锥内）
\end{itemize}

\subsection{摩擦锥约束}

每个$F_i$必须在相应的摩擦锥内：

\begin{itemize}
    \item 对于平面问题，每个摩擦锥由两条边定义
    \item $F_i$必须是摩擦锥边的非负线性组合
    \item 这可以表示为：$F_i = \sum_j \alpha_{ij} F_{ij}$，其中$\alpha_{ij} \geq 0$，$F_{ij}$是摩擦锥的边
\end{itemize}

\section{求解方法}

\subsection{线性规划}

这个线性约束满足问题可以通过多种方法求解，包括\textbf{线性规划}。

\subsubsection{线性规划形式}

可以将问题表述为线性规划问题：

\begin{align}
\text{寻找} \quad & k_i, \quad i = 1, \ldots, 16 \\
\text{最小化} \quad & \mathbf{1}^T \mathbf{k} \quad \text{（或任意目标函数）} \\
\text{使得} \quad & \text{方程(\ref{eq:system1})-(\ref{eq:system3})成立} \\
& k_i \geq 0, \quad i = 1, \ldots, 16 \\
& F_i \in \mathcal{WC}_i, \quad i = 1, \ldots, 16
\end{align}

\subsubsection{摩擦锥的线性化}

为了使用线性规划，需要将摩擦锥线性化：

\begin{itemize}
    \item 用多面体锥近似摩擦锥
    \item 每个摩擦锥由有限条边定义
    \item $F_i$必须是这些边的非负线性组合
\end{itemize}

\subsection{可行性测试}

\begin{itemize}
    \item \textbf{如果存在解}：装配在重力下是稳定的
    \item \textbf{如果不存在解}：装配不稳定，物体会移动或分离
\end{itemize}

\section{稳定性分析}

\subsection{稳定的条件}

装配在重力下稳定，如果：

\begin{enumerate}
    \item 存在满足所有平衡方程的接触力
    \item 所有接触力都在相应的摩擦锥内
    \item 所有接触系数$k_i \geq 0$
    \item 满足作用-反作用原理
\end{enumerate}

\subsection{不稳定的情况}

装配不稳定，如果：

\begin{enumerate}
    \item 不存在满足所有约束的解
    \item 某些接触必须分离（$B$模式）
    \item 某些接触必须滑动（$S$模式）
    \item 物体开始移动
\end{enumerate}

\section{扩展分析}

\subsection{其他接触模式}

虽然我们主要测试了所有接触都是$R$模式的情况，但也可以测试其他接触模式：

\begin{itemize}
    \item 某些接触滑动（$S$模式）
    \item 某些接触分离（$B$模式）
    \item 混合模式
\end{itemize}

\subsection{动态分析}

如果需要考虑动态效应：

\begin{itemize}
    \item 使用完整的动力学方程(12.31)
    \item 考虑加速度项
    \item 分析接触模式的转换
\end{itemize}

\section{关键要点总结}

\begin{enumerate}
    \item \textbf{多物体系统}：需要为每个物体建立平衡方程
    
    \item \textbf{作用-反作用原理}：接触力是成对出现的
    
    \item \textbf{代数方法}：对于复杂系统，代数方法比图形方法更适用
    
    \item \textbf{线性规划}：可以使用线性规划求解约束满足问题
    
    \item \textbf{摩擦锥约束}：每个接触力必须在摩擦锥内
    
    \item \textbf{可行性测试}：通过测试是否存在可行解来判断稳定性
\end{enumerate}

\section{实际应用}

这个例子展示了如何分析：

\begin{itemize}
    \item \textbf{装配稳定性}：判断装配是否会在重力下保持稳定
    \item \textbf{结构分析}：分析多物体系统的平衡
    \item \textbf{接触力计算}：计算保持平衡所需的接触力
    \item \textbf{设计验证}：验证设计是否满足稳定性要求
\end{itemize}

这些方法在以下应用中很重要：
\begin{itemize}
    \item 机器人装配
    \item 结构工程
    \item 产品设计
    \item 稳定性分析
\end{itemize}

\end{document}

